% !TeX root = ../../Book2.tex
\chapter{二阶偏微分方程的分类}
\label{b2:ch:08}
\BookChapterNumber{b2:ch:08}
\begin{chapterabstract}
二阶方程的主部在每一点确定一个实二次型。其正负方向与退化方向，把 Laplace、波动与热方程联系到三种不同的几何结构。本章定义主符号与特征余向量，推导坐标变换下的合同关系及低阶项，再以二维判别式和高维惯性说明经典分类的范围。二维常系数方程可完全化为标准型；变系数双曲和秩一方程可在适当邻域中构造特征坐标。椭圆主部的特殊系数变换与极坐标另行讨论，不能与实特征坐标的构造混同。时间导数的阶数、方向和系数符号将单独保留，以免把二阶矩阵的退化误当成完整的热型演化理论。
\end{chapterabstract}

\begin{learninggoals}
\begin{itemize}
\item\label{b2:item:08-goal-symbol}能正确提取对称主部矩阵和主符号，区分 Fourier 乘子、特征余向量与曲线切向。
\item\label{b2:item:08-goal-type}能判断二维类型与高维惯性，区分点态椭圆、一致椭圆、双曲时间方向和热型演化结构。
\item\label{b2:item:08-goal-coordinate}能逐项计算二阶坐标变换，包括由坐标二阶导数产生的一阶项，并给出合法的二维标准型；区分特征坐标、椭圆主部的特殊坐标以及极坐标各自适用的范围。
\item\label{b2:item:08-goal-limits}能识别退化线、非可积的候选坐标场及不稳定的数据面，说明分类能够提供什么、还需证明什么。
\end{itemize}
\end{learninggoals}

在 \(u_{xx}+u_{yy}=0\) 中把一个正号改成负号，就得到 \(u_{xx}-u_{yy}=0\)。二阶导数的阶数与数目没有改变，实特征线却从不存在变成两族直线。若换成 \(u_y-u_{xx}=0\)，又出现只沿一个方向演化的扩散。要判断应该使用边界约束、传播还是时间耗散，首先需要从主部的二次型读出其几何结构，同时保留时间的一阶项。

\ProseParagraphBreak

本章使用有限维实对称矩阵的谱定理、\chapref{b2:ch:01}的二阶链式法则与逆函数定理，以及\chapref{b2:ch:07}的局部流。\secref{b2:sec:0801}介绍坐标无关的分类概念；\secref{b2:sec:0802}先用\cref{b2:tab:08-planar-classification}集中比较二维判别式与实特征方向，再发展高维类型及时间方向。\secref{b2:sec:0803}从\subsecref{b2:subsec:080301}的标准型和坐标导数可积性出发，构造变系数特征坐标；\subsecref{b2:subsec:08-elliptic-polar-coordinates}单独讨论椭圆主部的特殊坐标与极坐标。初读时可先掌握二维判别式和三种模型，再回看高维惯性与时间法向的区别。

\ProseParagraphBreak

Walter A. Strauss 的《Partial Differential Equations: An Introduction》\cite[Section 1.6]{Strauss2008Partial}用实二次型讲解二维标准型与高维分类。其按惯性给出的抛物型名称在本章保留，但本章另外写清时间一阶项的作用，避免只凭主部秩就推断初值问题的性质。变系数坐标构造使用前章建立的局部流，不要求先掌握后续波动或热方程的解公式。
% Strauss 第二版 Wiley 官方样章印刷28--31页已实读；常系数标准型为独立组织的完整推导。

% !TeX root = ../../Book2.tex
\section{二阶线性方程}
\label{b2:sec:0801}

两个二阶方程可能只差一个符号，解的行为却截然不同。分类首先抓住最高阶导数的组合，把它看成每一点上的二次型。本节使用 \(z=(z_1,\ldots,z_n)\) 表示全部自变量；只有在讨论演化问题时才拆成 \(z=(t,x)\)。因此 \(n\) 不总等于空间维数。

\subsection{主部}
\label{b2:subsec:080101}

在开集 \(\Omega\subseteq\mathbb R^n\) 上考虑实系数算子
\begin{equation}\label{b2:eq:08-linear-operator}
 Lu=\sum_{i,j=1}^n a^{ij}(z)\partial_i\partial_ju
 +\sum_{i=1}^n b^i(z)\partial_i u+c(z)u.
\end{equation}
系数至少连续，经典解至少属于 \(C^2\)。一阶和零阶项称为低阶项；主部已在\secref{b2:sec:0104}引入，此处专指含二阶导数的和式，记为 \(L_2u\)。由于混合偏导相等，把 \(a^{ij}\) 换成 \((a^{ij}+a^{ji})/2\) 不改变算子。以下始终取实对称矩阵 \(A=(a^{ij})\)。

\ProseParagraphBreak

二维中习惯写为
\[
 L_2u=a u_{xx}+2b u_{xy}+c u_{yy},\quad
 A=\begin{pmatrix}a&b\\b&c\end{pmatrix}.
\]
混合项前面的 \(2\) 来自矩阵求和中 \((1,2)\) 和 \((2,1)\) 两项。若方程写成 \(u_{xx}+6u_{xy}+u_{yy}=0\)，则矩阵的非对角元为 \(3\)，不是 \(6\)。本节中的 \(a,b,c\) 是二维主部的临时系数，不与\eqref{b2:eq:08-linear-operator}中的向量低阶系数混用。

\ProseParagraphBreak

散度形式也有相同的最高阶矩阵。在 \(a^{ij}\in C^1\) 时，
\[
 \sum_{i,j}\partial_i(a^{ij}\partial_ju)
 =\sum_{i,j}a^{ij}\partial_i\partial_ju
 +\sum_{i,j}(\partial_i a^{ij})\partial_ju.
\]
它与非散度形式在经典分类上相同，但分部积分时低阶项不同；在粗糙系数或弱解问题中，两种写法更不能只凭主部相同便随意交换。

\subsection{主符号}
\label{b2:subsec:080102}

\begin{definition}\label{b2:def:08-principal-symbol}
设 \(n\geq1\)，\(\Omega\subset\mathbb R^n\) 为开集，实系数算子 \(L=\sum_{i,j}a^{ij}(z)\partial_{ij}+\sum_i b^i(z)\partial_i+c(z)\) 的系数均连续，且 \(A=(a^{ij})\) 为实对称矩阵场。对 \(z\in\Omega\)，二次齐次多项式
\[
 p(z,\xi)\coloneqq\xi^{\mathsf T}A(z)\xi
 =\sum_{i,j=1}^n a^{ij}(z)\xi_i\xi_j
 \quad(\xi\in\mathbb R^n)
\]
称为 \(L\) 的\term[zhufuhao]{主符号}[principal symbol]。本章不把 Fourier 微分规则中的 \(-4\pi^2\) 因子纳入这个几何记号。
\end{definition}
\symidx{\(p(z,\xi)\)：二阶算子的几何主符号}

后一约定与第二编的 \(2\pi\) 归一化相容：对常系数主部，
\[
 L_2e^{2\pi\mathrm i z\cdot\xi}
 =-4\pi^2p(\xi)\cdot e^{2\pi\mathrm i z\cdot\xi}.
\]
因此 \(p\) 记录不同频率方向的二阶作用，整个 Fourier 乘子则是 \(-4\pi^2p\)。不同文献可能把这个负号纳入符号定义；判断椭圆性和特征集合时，必须先核对其约定。

\ProseParagraphBreak

更一般地，对光滑实相位 \(\phi\) 与大参数 \(\lambda\)，直接两次求导给出
\[
 L\bigl(e^{2\pi\mathrm i\lambda\phi}\bigr)
 =e^{2\pi\mathrm i\lambda\phi}
 \left(-4\pi^2\lambda^2p(z,\nabla\phi)
 +2\pi\mathrm i\lambda\left(\sum_{i,j}a^{ij}\partial_{ij}\phi
 +\sum_i b^i\partial_i\phi\right)+c\right).
\]
在固定紧集上，最高的 \(\lambda^2\) 次项由主符号决定。若该系数为零，通常占主导的二阶作用恰好抵消；特征几何正由这类方向产生。这个计算只是代数恒等式，没有断言任意高频函数就是某个方程的近似解。

\subsection{特征余向量}
\label{b2:subsec:080103}

\begin{definition}\label{b2:def:08-characteristic-covector}
设 \(n\geq1\)，\(\Omega\subset\mathbb R^n\) 为开集，实系数算子 \(L=\sum_{i,j}a^{ij}(z)\partial_{ij}+\sum_i b^i(z)\partial_i+c(z)\) 的系数均连续，且 \(A=(a^{ij})\) 为实对称矩阵场。令 \(p(z,\nu)=\nu^{\mathsf T}A(z)\nu\)。在一点 \(z\in\Omega\)，若非零实余向量 \(\nu\in(\mathbb R^n)^*\) 满足 \(p(z,\nu)=0\)，则称它为该点的\term[tezheng-yuxiangliang]{特征余向量}[characteristic covector]。若实值函数 \(\phi\in C^2(V)\) 定义在开集 \(V\subseteq\Omega\) 上，在超曲面 \(\Sigma=\{z\in V:\phi(z)=0\}\) 上处处有 \(\nabla\phi\ne0\)，且 \(p(z,\mathrm d\phi(z))=0\)，则称 \(\Sigma\) 为\term[tezheng-chaqumian]{特征超曲面}[characteristic hypersurface]。在 \(\Sigma\) 上满足 \(p(z,\mathrm d\phi(z))\ne0\) 的点称为该曲面的非特征点。
\end{definition}

这里的 \(\nu\) 表示余向量，也就是对位移作线性求值的法向数据；借欧氏内积可用列向量表示，但不能把它直接当作曲面的切向。例如，\(\mathrm d\phi(z)\) 对切向量 \(w\) 的作用是 \(\mathrm d\phi(z)[w]=\nabla\phi(z)\cdot w\)，在切平面上等于零。下文计算主符号时，以梯度的列向量表示同一个微分。非零倍数 \(\alpha\nu\) 仍为特征余向量，所以也可以讨论它确定的、忽略非零实倍数的余法向方向；二维特征曲线的切方向则是另一个对象。

\ProseParagraphBreak

换一个定义函数 \(\widetilde\phi=\alpha\phi\)，其中 \(\alpha\) 是曲面附近不为零的 \(C^2\) 实函数，在曲面上有 \(\mathrm d\widetilde\phi=\alpha\,\mathrm d\phi\)。主符号齐次，零值条件不变。一般地，同一光滑超曲面的两个正则定义函数在曲面上的非零微分都有同一个核，即切平面，因而成非零比例；所以特征条件不依赖定义函数的选择。一个曲面也可能只有部分点为特征点，不能从一个点的条件推出整张曲面均为特征。

\ProseParagraphBreak

在二维中，若曲线切向为 \((\mathrm dx,\mathrm dy)\)，则其一个法向与 \((-\mathrm dy,\mathrm dx)\) 平行。特征条件因而写作
\begin{equation}\label{b2:eq:08-characteristic-differential}
 a(\mathrm dy)^2-2b\,\mathrm dx\,\mathrm dy+c(\mathrm dx)^2=0.
\end{equation}
当 \(\mathrm dx\ne0\) 时，它变为斜率方程 \(am^2-2bm+c=0\)。这解释了为什么这里混合项的符号与主符号 \(a\xi^2+2b\xi\eta+c\eta^2\) 相反：一个写切向，一个写法向。

\ProseParagraphBreak

非特征条件的代数意义是可以解出法向二阶导数。取坐标 \(y_1=\phi(z)\)，补足其余坐标后，\(\partial_{y_1}^2v\) 的系数就是 \(p(z,\nabla\phi)\)。该系数非零时，可以把它移到方程一侧并相除。但这一步只整理了方程，没有证明给定 Cauchy 数据的解存在，也没有得到稳定性。例如对 Laplace 方程，任意非零实法向都是非特征的，而给定函数值和法向导数的 Cauchy 问题仍可能非常不稳定。

\subsection{坐标不变性}
\label{b2:subsec:080104}

\begin{proposition}\label{b2:prop:08-coordinate-change}
设 \(n\geq1\)，\(\Omega,\widetilde\Omega\subset\mathbb R^n\) 为开集，实系数算子 \(L=\sum_{i,j}a^{ij}(z)\partial_{ij}+\sum_i b^i(z)\partial_i+c(z)\) 的系数均连续，且 \(A=(a^{ij})\) 为实对称矩阵场。设 \(F:\Omega\to\widetilde\Omega\) 为 \(C^2\) 微分同胚，\(\mathbb K\in\{\mathbb R,\mathbb C\}\)，\(v\in C^2(\widetilde\Omega;\mathbb K)\)，\(y=F(z)\)，\(u(z)=v(F(z))\)，并令 \(B=DF(z)\)。则新坐标中的主部矩阵为
\begin{equation}\label{b2:eq:08-congruence}
 \widetilde A(y)=B A(z)B^{\mathsf T},\quad z=F^{-1}(y),
\end{equation}
一阶系数为
\begin{equation}\label{b2:eq:08-lower-transform}
 \widetilde b^k(y)=\sum_{i,j}a^{ij}(z)\partial_{ij}F^k(z)
 +\sum_i b^i(z)\partial_iF^k(z).
\end{equation}
零阶系数为 \(\widetilde c(y)=c(F^{-1}(y))\)，主符号满足 \(\widetilde p(y,\eta)=p(z,B^{\mathsf T}\eta)\)。因此，对 \(f\in C(\Omega;\mathbb K)\)，方程 \(Lu=f\) 等价于
\[
 \sum_{k,l}\widetilde a^{kl}\partial_{kl}v
 +\sum_k\widetilde b^k\partial_kv+\widetilde c\,v
 =\widetilde f,
 \qquad \widetilde f=f\circ F^{-1}.
\]
\end{proposition}
\begin{proof}
链式法则给出
\[
 \partial_i u=\sum_k\partial_kv\cdot\partial_iF^k,
\quad
 \partial_{ij}u=\sum_{k,l}\partial_{kl}v\cdot\partial_iF^k\cdot\partial_jF^l
 +\sum_k\partial_kv\cdot\partial_{ij}F^k.
\]
代入算子并按 \(v\) 的导数阶数整理，得到两条系数公式。再把 \(\eta\) 代入二次型，便有
\(\eta^{\mathsf T}BAB^{\mathsf T}\eta=(B^{\mathsf T}\eta)^{\mathsf T}A(B^{\mathsf T}\eta)\)。
\end{proof}

实对称矩阵在可逆\term[hetong-bianhuan]{合同变换}[congruence transformation]下保持正、负及零特征值的个数，这三项称为其\term[guanxing]{惯性}[inertia]。为说明这一事实，先由有限维谱定理把矩阵正交对角化。二次型为正定的子空间最大维数恰等于正特征值个数：正特征子空间达到这个维数；任何更大子空间必与非正特征子空间非平凡相交，不可能正定。可逆线性变换保持子空间维数，因而保持该最大值。对负号二次型重复此论证可得负特征值个数不变，最后由总维数得到零特征值个数不变。

\ProseParagraphBreak

所以主部分类不依赖合法坐标，但矩阵的具体特征值和统一常数可能改变。设 \(A(z)\geq\lambda I\)，其中 \(\lambda>0\) 与 \(z\) 无关，则
\[
 \eta^{\mathsf T}\widetilde A(F(z))\eta
 \geq\lambda|B(z)^{\mathsf T}\eta|^2.
\]
若还知道 \(\sup_{z\in\Omega}\|B(z)^{-1}\|\leq M<\infty\)，由
\( |\eta|\leq\|B(z)^{-1}\|\,|B(z)^{\mathsf T}\eta|\)
得到 \(\widetilde A\geq\lambda M^{-2}I\)。这里均使用欧氏算子范数，且转置不改变该范数。逆导数的一致控制是从 \(A\geq\lambda I\) 这一已知下界推出新下界的充分条件；缺少它时，这一个下界本身不足以保证结论，仍须检查 \(A\) 与 \(B\) 的共同变化。

\ProseParagraphBreak

例如在 \(\mathbb R^2\) 上取
\[
 A(x,y)=\begin{pmatrix}1&0\\0&(1+y^2)^2\end{pmatrix},
 \qquad F(x,y)=(x,\arctan y).
\]
这是到 \(\mathbb R\times(-\pi/2,\pi/2)\) 的光滑微分同胚，且
\[
 B=\begin{pmatrix}1&0\\0&(1+y^2)^{-1}\end{pmatrix},
 \qquad \|B^{-1}\|=1+y^2,
 \qquad BAB^{\mathsf T}=I.
\]
逆导数虽然无界，原矩阵在同一方向上的增长却恰好抵消坐标压缩，新矩阵仍有统一正下界。反过来，若保持同一个 \(F\) 而取 \(A=I\)，新矩阵第二个特征值为 \((1+y^2)^{-2}\)，确实没有统一正下界。这两个计算同时说明充分条件的用途和它的适用边界。

\ProseParagraphBreak

还须区分统一正下界与\secref{b2:sec:0802}采用的双侧一致椭圆条件。前一个例子的原矩阵没有统一上界，不能称为该双侧意义下的一致椭圆矩阵。若已知 \(\lambda I\leq A\leq\Lambda I\)，再有 \(\sup\|B\|\leq N\)、\(\sup\|B^{-1}\|\leq M\)，则
\[
 \frac{\lambda}{M^2}|\eta|^2
 \leq\eta^{\mathsf T}\widetilde A\eta
 \leq\Lambda N^2|\eta|^2,
\]
从而双侧条件也被传递。点态坐标不变性与全域定量估计是两个层次。非线性坐标还会产生\eqref{b2:eq:08-lower-transform}中的一阶项，忽略它们会改变原方程。

% !TeX root = ../../Book2.tex
\section{三种基本类型}
\label{b2:sec:0802}

本节先用二维判别式建立直观，再按实对称主部的惯性给出高维定义，最后说明时间方向及低阶项带来的额外区别。分类针对实系数标量算子，未知函数可以取实值或复值；若主系数本身为复数，或考虑方程组、高阶方程，则需另定分类方法。

\ProseParagraphBreak

先在一点固定二维主部 \(a u_{xx}+2b u_{xy}+c u_{yy}\)，其中 \(a,b,c\in\mathbb R\)，记
\[
 A=\begin{pmatrix}a&b\\b&c\end{pmatrix},\quad D=b^2-ac=-\det A.
\]
\Cref{b2:tab:08-planar-classification}中的一个“方向”代表一条过原点的直线，不区分互为非零实倍数的余向量，特别是不区分 \(\nu\) 与 \(-\nu\)。特征余向量满足 \(\nu^{\mathsf T}A\nu=0\)；对应曲线的切方向是它的核直线，两者不能混作同一个向量。

\begin{table}[htbp]
\centering
\small
\setlength{\tabcolsep}{4pt}
\caption[二维主部的代数分类]{二维主部的代数分类。\(A\) 为实对称 \(2\times2\) 矩阵，\(D=b^2-ac\)；只计某一点的非定向实特征方向。常系数时，不同切方向分别给出平行特征直线族。}
\label{b2:tab:08-planar-classification}
\begin{tabularx}{\linewidth}{@{}>{\raggedright\arraybackslash}p{2.3cm}>{\raggedright\arraybackslash}p{3.8cm}>{\raggedright\arraybackslash}p{3.2cm}>{\raggedright\arraybackslash}X@{}}
\toprule
条件 & 惯性及代数类型 & 特征余方向 & 特征曲线切方向 \\
\midrule
\(D<0\) & 正定或负定；椭圆型 & 无 & 无 \\
\(D>0\) & 一正一负；双曲型 & 两个不同方向 & 两个不同方向 \\
\(D=0,\ A\ne0\) & 秩一；抛物型的代数类型 & 一个方向，即 \(\ker A\) & 一个方向 \\
\(A=0\) & 秩零；二阶主部全退化 & 所有非零余向量均使主符号为零 & 所有切方向均满足零符号条件 \\
\bottomrule
\end{tabularx}
\end{table}

这个表可由实对称矩阵的正交对角化直接核对：正定或负定二次型没有非零实零向量；异号二次型的零集为两条不同直线；秩一半定二次型的零集只有核直线。每个余方向又唯一确定一个垂直的切方向。于是三种类型覆盖所有非零的二维实二次型；\(A=0\) 必须单列，不能仅因 \(D=0\) 就叫作通常的抛物型。表内“抛物型”只说二阶主部，热型初值问题的演化方向与适定性还需要低阶项及数据条件，后文将具体比较。

\subsection{椭圆型}
\label{b2:subsec:080201}

\begin{definition}\label{b2:def:08-elliptic}
设 \(n\geq1\) 为整数，\(\Omega\subset\mathbb R^n\) 为开集，\(L=\sum_{i,j=1}^n a^{ij}(z)\partial_{ij}+\sum_{i=1}^n b^i(z)\partial_i+c(z)\)，所有系数为实值连续函数，\(A=(a^{ij})\) 为实对称矩阵场。若 \(A(z)\) 正定或负定，则称二阶算子在 \(z\in\Omega\) 为\term[tuoyuanxing]{椭圆型}[elliptic]。若它在开集每一点均为椭圆型，就称在该开集为椭圆型。选定正定的符号后，若存在常数 \(0<\lambda\leq\Lambda<\infty\)，使
\[
 \lambda|\xi|^2\leq\xi^{\mathsf T}A(z)\xi\leq\Lambda|\xi|^2
 \quad(z\in\Omega,\ \xi\in\mathbb R^n),
\]
则称主部为\term[yizhi-tuoyuan]{一致椭圆}[uniformly elliptic]。
\end{definition}

对实对称矩阵，椭圆性也等价于对每个 \(\xi\ne0\) 都有 \(p(z,\xi)\ne0\)。若特征值中含零，存在非零零向量；若有正、负特征值，取相应单位特征向量 \(e_+,e_-\)，二次型在 \(e_+\cos\theta+e_-\sin\theta\) 上从正变负，也必取零。反过来，正定或负定直接排除零值。因此椭圆主部没有实特征超曲面。

\ProseParagraphBreak

最基本的模型为 \(\Delta u=f\) 或 \(-\Delta u=f\)。对称正定但随位置退化的系数提醒我们，点态椭圆不等于一致椭圆：在半平面 \(x>0\) 上，\(u_{xx}+x^2u_{yy}\) 处处椭圆，但沿 \(x\downarrow0\) 的最小特征值趋零，沿 \(x\to\infty\) 的最大特征值无界。在任意紧子集上，这两种障碍都消失。

\begin{proposition}\label{b2:prop:08-compact-ellipticity}
设 \(n\geq1\) 为整数，\(\Omega\subset\mathbb R^n\) 为连通开集，\(A:\Omega\to\mathbb R^{n\times n}\) 连续且取实对称矩阵值，在 \(\Omega\) 中每一点均正定或负定。则它的正、负定符号在整个开集上一致。选为正定后，对每个紧集 \(K\subseteq\Omega\)，都有依赖于 \(K\) 的一致椭圆常数。
\end{proposition}
\begin{proof}
正定矩阵构成开集：若 \(\min_{|\xi|=1}\xi^{\mathsf T}A(z_0)\xi=\lambda_0>0\)，则 \(\|A(z)-A(z_0)\|<\lambda_0/2\) 时仍正定。负定情形相同。两类点各为开集且不交，连通性迫使只出现一类。选好符号后，若 \(K\ne\varnothing\)，连续函数 \((z,\xi)\mapsto\xi^{\mathsf T}A(z)\xi\) 在紧集 \(K\times S^{n-1}\) 上取得正最小值及有限最大值，即为所需常数；空集上的结论直接成立。
\end{proof}

\subsection{双曲型}
\label{b2:subsec:080202}

\begin{definition}\label{b2:def:08-hyperbolic}
设 \(n\geq2\) 为整数，\(\Omega\subset\mathbb R^n\) 为开集，\(L=\sum_{i,j=1}^n a^{ij}(z)\partial_{ij}+\sum_{i=1}^n b^i(z)\partial_i+c(z)\)，所有系数为实值连续函数，\(A=(a^{ij})\) 为实对称矩阵场。若 \(A(z)\) 非退化，且一个特征值与其余 \(n-1\) 个特征值符号相反，则称算子在 \(z\in\Omega\) 为\term[shuangquxing]{双曲型}[hyperbolic]。在演化问题中，还要选定时间余法向，并核对相对于它的时间多项式，而不只检查法向非特征。
\end{definition}

设 \(d\geq1\) 为整数、\(c>0\)。波动算子 \(\partial_t^2-c^2\Delta_x\) 作用于 \(d\) 个空间变量和一个时间变量，其矩阵惯性为 \((1,d,0)\)，主符号为
\[
 p(\tau,\xi)=\tau^2-c^2|\xi|^2.
\]
这里 \(\xi\in\mathbb R^d\) 是空间频率，\(\tau\) 是时间频率。若只看空间块 \(-c^2I\)，它是负定的；整个时空主部才是双曲型。称“波动方程的空间部分椭圆”与称“波动方程双曲”并不矛盾。

\ProseParagraphBreak

“选择时间”有一个可计算判据。在一点固定符号使惯性为 \((1,n-1,0)\)，选择满足 \(p(\nu)>0\) 的法向 \(\nu\)。补足一组坐标，使 \(\nu\) 对应时间余方向，则矩阵分块为
\[
 A=\begin{pmatrix}a&v^{\mathsf T}\\v&C\end{pmatrix},\quad a>0.
\]
配方给出
\begin{equation}\label{b2:eq:08-hyperbolic-square}
 p(\tau,\xi)=a\left(\tau+\frac{v\cdot\xi}{a}\right)^2
 +\xi^{\mathsf T}\left(C-\frac{vv^{\mathsf T}}{a}\right)\xi.
\end{equation}
右侧的空间矩阵必须负定，因为合同消元后已有一个正方向，全部剩余方向均为负。故每个非零空间频率对应两个不同的实根
\[
 \tau=-\frac{v\cdot\xi}{a}
 \pm\sqrt{-\frac1a\xi^{\mathsf T}\left(C-\frac{vv^{\mathsf T}}{a}\right)\xi}.
\]
这种相对于时间余法向的实根性质称为\term[yange-shuangqu]{严格双曲}[strictly hyperbolic]。精确地说，固定一个非零余法向 \(\nu\)，把它补足为一组余空间基，后 \(n-1\) 个坐标称为空间余方向；要求 \(p(\nu)\ne0\)，并且对每个非零空间频率 \(\xi\in\mathbb R^{n-1}\)，关于 \(\tau\) 的二次多项式 \(p(\tau,\xi)\) 都有两个不同的实根。改变空间补足坐标只会重标空间频率并平移时间根，不改变这一性质。上面的配方证明 \(p(\nu)>0\) 是所选惯性 \((1,n-1,0)\) 下的充分条件；这一判据针对二阶标量方程的两个不同实特征根。

\ProseParagraphBreak

在至少两个空间变量的波动模型中，任意非特征余法向未必都能作为这样的时间方向。对三变量模型 \(p(\tau,\xi_1,\xi_2)=\tau^2-\xi_1^2-\xi_2^2\)，若把 \(x_1\) 当时间，那么固定原来的 \(\tau=0\)、\(\xi_2\ne0\) 时，关于 \(\xi_1\) 的方程没有实根。这里原惯性中的唯一正方向与其余负方向有不同地位；相对于选定数据面的实根性质须另行检查。

\ProseParagraphBreak

倾斜的时间面给出一个定量例子。对 \(d\geq2\) 个空间变量的波动算子 \(\partial_t^2-c^2\Delta_x\)，其中 \(c>0\)，取 \(T=t+\kappa x_1\)、\(X_j=x_j\)（\(1\leq j\leq d\)），\(\kappa\in\mathbb R\)。逆变换为 \(t=T-\kappa X_1\)，所以对所有 \(\kappa\) 都可逆。由 \(\partial_t=\partial_T\)、\(\partial_{x_1}=\kappa\partial_T+\partial_{X_1}\)，新主部为
\[
 (1-c^2\kappa^2)\partial_T^2
 -2c^2\kappa\partial_T\partial_{X_1}
 -c^2\sum_{j=1}^d\partial_{X_j}^2.
\]
余法向 \(\mathrm{d}T=\mathrm{d}t+\kappa\,\mathrm{d}x_1\) 的主符号为 \(p(\mathrm{d}T)=1-c^2\kappa^2\)，记此值为 \(a\)。令 \(\eta=(\eta_1,\eta_\perp)\in\mathbb R\times\mathbb R^{d-1}\)，则时间多项式及其判别式分别为
\[
 p_T(\tau,\eta)=a\tau^2-2c^2\kappa\eta_1\tau-c^2|\eta|^2,
 \quad \Delta_\tau=4c^2\bigl(\eta_1^2+a|\eta_\perp|^2\bigr).
\]
若 \(|\kappa|<1/c\)，则 \(a>0\)，每个 \(\eta\ne0\) 对应的两个时间根为
\[
 \tau=\frac{c^2\kappa\eta_1\pm c\sqrt{\eta_1^2+a|\eta_\perp|^2}}a,
\]
它们均实且互异。若 \(|\kappa|=1/c\)，等 \(T\) 面是特征面，最高时间导数的系数消失。若 \(|\kappa|>1/c\)，它虽非特征，但可取 \(\eta_1=0\)、\(\eta_\perp\ne0\)，此时 \(a\tau^2-c^2|\eta_\perp|^2=0\) 没有实根。因此在 \(d\geq2\) 时，只有 \(|\kappa|<1/c\) 才使这一倾斜时间方向严格双曲。整个矩阵的惯性一直不变，所检查的是数据面相对于特征锥的实根性质。

\ProseParagraphBreak

这里必须保留维数条件。若只有一个空间变量，即惯性为 \((1,1,0)\)，便没有 \(\eta_\perp\)；对任意 \(a\ne0\) 及 \(\eta_1\ne0\)，判别式均为 \(4c^2\eta_1^2>0\)。更一般地，对任意惯性为 \((1,1,0)\) 的实对称矩阵 \(A\)，取非特征余向量 \(\nu\) 及与它线性独立的 \(\zeta\)。多项式 \(p(\tau\nu+\zeta)\) 的判别式为
\[
 4\bigl((\nu^{\mathsf T}A\zeta)^2-p(\nu)p(\zeta)\bigr)
 =-4\det A\cdot\det(\nu,\zeta)^2>0.
\]
等式来自基矩阵 \((\nu,\zeta)\) 下二次型的行列式公式。因此 \(1+1\) 维中任一非特征法向都给出严格双曲性，包括 \(p(\nu)<0\) 的情形；整体乘以负号后，就能对后者使用时间系数为正的配方。若预先坚持选用包含通常 \(\mathrm{d}t\) 的正号时间锥，则此倾斜族仍取 \(a>0\)，但这项时间锥选择须与“是否满足实根判据”区分。

\Cref{b2:fig:08-tilted-time}用相同的 \((x_1/c,t)\) 尺度画出三种倾斜数据面，并标记余法向主符号的符号。图中的 \(1+1\) 维截面用于辨认角度及特征边界；\(d\geq2\) 时负号法向的失败，还须使用刚才的横向频率 \(\eta_\perp\)，不能只从这张截面图推出。
% 本书原创：统一使用横轴 x_1/c、纵轴 t，比较三种倾斜时间面的余法向符号。
\begin{figure}[htbp]
\centering
\begin{tikzpicture}[x=.9cm,y=.9cm,>=stealth,every node/.style={font=\small},line width=.65pt]
\foreach \offset/\slope/\ktxt/\psign in {
  0/.5/{\kappa=\dfrac{1}{2c}}/{p(\mathrm{d}T)>0},
  4.4/1/{\kappa=\dfrac1c}/{p(\mathrm{d}T)=0},
  8.8/1.5/{\kappa=\dfrac{3}{2c}}/{p(\mathrm{d}T)<0}}{
 \begin{scope}[shift={(\offset,0)}]
  \begin{scope}
   \clip (-1.25,-1.25) rectangle (1.25,1.25);
   \draw[gray,dashed] (-1.25,-1.25)--(1.25,1.25);
   \draw[gray,dashed] (-1.25,1.25)--(1.25,-1.25);
   \draw[AnalysisBlue,very thick,domain=-1.25:1.25,samples=2,variable=\q]
     plot ({\q},{-\slope*\q});
  \end{scope}
  \draw[->,gray] (-1.4,0)--(1.45,0) node[below right] {\(x_1/c\)};
  \draw[->,gray] (0,-1.4)--(0,1.45) node[above] {\(t\)};
  \draw[->,orange!70!black,thick] (0,0)--({.6*\slope},.6)
     node[above right,fill=white,inner sep=1pt] {\(\mathrm{d}T\)};
  \node at (0,2.5) {\(\ktxt\)};
  \node at (0,-1.8) {\(\psign\)};
 \end{scope}
}
\node[align=center,text width=12cm] at (4.4,-2.65)
 {蓝线：等 \(T\) 数据面 \(T=0\)；灰虚线：\(t=\pm x_1/c\)；\par
 箭头：余法向 \(\mathrm{d}T\) 的欧氏表示。};
\end{tikzpicture}
\caption[倾斜时间面与特征线]{倾斜时间面与特征线。波动算子为 \(\partial_t^2-c^2\Delta_x\)，\(c>0\)，\(T=t+\kappa x_1\)。三图采用相同的横纵尺度；\(\kappa<0\) 时左右反射。这里是 \(1+1\) 维截面，右图的非特征法向在 \(1+1\) 维仍严格双曲；至少两个空间变量时的失败还须使用正文中的横向频率。}
\label{b2:fig:08-tilted-time}
\end{figure}


\subsection{抛物型}
\label{b2:subsec:080203}

\begin{definition}\label{b2:def:08-parabolic-algebraic}
设 \(n\geq2\) 为整数，\(\Omega\subset\mathbb R^n\) 为开集，\(L=\sum_{i,j=1}^n a^{ij}(z)\partial_{ij}+\sum_{i=1}^n b^i(z)\partial_i+c(z)\)，所有系数为实值连续函数，\(A=(a^{ij})\) 为实对称矩阵场。若 \(A(z)\) 的秩为 \(n-1\)，且所有非零特征值同号，则称它的主部在 \(z\in\Omega\) 具有\term[paowuxing]{抛物型}[parabolic]的代数类型。这个条件只描述二阶部分，不单独保证方程是一条随时间演化的热型方程。
\end{definition}

例如对 \(d\geq1\) 个空间变量，\(\partial_t-\Delta_x\) 对全部 \(d+1\) 个变量的二阶矩阵为 \(\operatorname{diag}(0,-1,\ldots,-1)\)。时间余法向 \((1,0)\) 使二阶主符号为零，所以等时面在刚才的纯二阶定义下是特征的。热方程仍可给初值，因为沿时间方向真正出现的是一阶导数。把所有变量都按同一阶数计算的主符号，不能独自表达这种演化结构。

\ProseParagraphBreak

设 \(d\geq1\) 为整数，\(U\subset\mathbb R^d\times\mathbb R\) 为开集。对 \(U\) 上的实值连续系数及连续右端 \(f\)，本编的热型演化模型特指
\begin{equation}\label{b2:eq:08-parabolic-evolution}
 u_t-\sum_{i,j=1}^d a^{ij}(x,t)u_{x_ix_j}
 +\sum_i b^i(x,t)u_{x_i}+c(x,t)u=f,
\end{equation}
其中空间矩阵 \((a^{ij})\) 对称正定；若存在常数 \(0<\lambda\leq\Lambda<\infty\)，使
\[
 \lambda|\xi|^2\leq\sum_{i,j=1}^d a^{ij}(x,t)\xi_i\xi_j
 \leq\Lambda|\xi|^2
 \quad((x,t)\in U,\ \xi\in\mathbb R^d),
\]
则称这一演化结构为一致抛物。时间一阶系数已规范为 \(1\)，其符号与空间二阶项相反。这项结构同时使用了一阶项和选定的时间方向。若时间导数前有正函数 \(\alpha\)，须先将整个方程除以 \(\alpha\)，再检查空间矩阵的正定性与统一常数。

\ProseParagraphBreak

以下三个算子有相同的二阶矩阵，却代表不同问题：
\[
 \partial_t-\partial_x^2,\quad
 -\partial_t-\partial_x^2,\quad
 -\partial_x^2.
\]
第一个是向前热算子；第二个反转时间后才成为向前热算子；第三个根本没有时间演化，只是每个 \(t\) 上的一维方程。由此不能仅凭“一个特征值为零”就声称存在通常的热方程初值理论。在一般高维变量中，秩还可能低于 \(n-1\)，或者非零特征值有正有负；这些属于其他退化类型。

\subsection{标准模型}
\label{b2:subsec:080204}

对二维非零实主部，\cref{b2:tab:08-planar-classification}已经用 \(D=b^2-ac\) 区分三个类型；表中的零矩阵情形则完全失去二阶主部。下面的模型说明，点态代数分类之外还须观察系数变化与演化方向。

\ProseParagraphBreak

分类可以随位置改变。以 \(Lu=y u_{xx}+u_{yy}\) 为例，当 \(y>0\) 时矩阵正定，当 \(y<0\) 时特征值异号，在 \(y=0\) 时秩为一。穿过退化线的一个邻域上不可能存在把主部变成某个固定非退化类型的微分同胚，因为合同变换保持各点的秩与惯性。退化线须作为独立几何对象处理，不能在变换时相除消去。

\ProseParagraphBreak

频率模型说明时间符号为何值得单独保留。在单位圆周 \(\mathbb T=\mathbb R/\mathbb Z\) 上使用总质量为 \(1\) 的 Lebesgue 测度，取 \(c>0\)、空间频率 \(k\in\mathbb Z\)，令 \(u(x,t)=v_k(t)e^{2\pi\mathrm i kx}\)，三种方程分别给出
\[
 \begin{aligned}
 u_{tt}-c^2u_{xx}=0&\ \Longrightarrow\ v_k''+4\pi^2c^2k^2v_k=0,\\
 u_t-u_{xx}=0&\ \Longrightarrow\ v_k'+4\pi^2k^2v_k=0,\\
 u_t+u_{xx}=0&\ \Longrightarrow\ v_k'-4\pi^2k^2v_k=0.
 \end{aligned}
\]
第一行在 \(k\ne0\) 时给出振荡，第二行使频率按 \(e^{-4\pi^2k^2t}\) 衰减，第三行则按 \(e^{4\pi^2k^2t}\) 放大。这些都是直接代入得到的确切解。对第三行取正整数 \(k\to\infty\) 及初值 \(u_k(x,0)=e^{-k}e^{2\pi\mathrm i kx}\)。对每个固定非负整数 \(m\)，令 \(\|h\|_{C^m(\mathbb T)}=\max_{0\leq j\leq m}\|\partial_x^jh\|_\infty\)，则
\[
 \|u_k(\cdot,0)\|_{C^m(\mathbb T)}
 \leq C_m k^m e^{-k}\longrightarrow0,
 \quad \|u_k(\cdot,t)\|_{L^2(\mathbb T)}
 =e^{4\pi^2k^2t-k}\longrightarrow\infty\quad(t>0).
\]
因此即使二阶代数类型相同，对同一时间方向的数据也可能有完全不同的稳定性。

\ProseParagraphBreak

低阶项也会影响唯一性、增长和边界相容性，\chapref{b2:ch:01}的能量例子已经显示这一点。这些类型提示了不同的分析方法：双曲型需要研究传播和能量，热型演化需要研究扩散和单向时间，椭圆型需要研究边界约束与内部性质。区域、边界数据和系数的正则性进一步决定各项估计的适用范围。

% !TeX root = ../../Book2.tex
\section{坐标变换}
\label{b2:sec:0803}

主符号的惯性告诉我们可以期待哪一种模型，坐标变换则把这一判断落实为计算。本节区分三种任务：常系数主部在整个坐标域中的标准化，变系数主部在一点的标准化，以及利用特征曲线在一个邻域中构造坐标。三者使用的条件不同，结论也不相同。

\subsection{二维标准型与坐标可积性}
\label{b2:subsec:080301}

\begin{theorem}\label{b2:thm:08-constant-canonical}
设 \(a,b,c\) 为实常数且不全为零。存在可逆实线性坐标变换，并允许把整个方程乘以一个非零常数，使主部 \(a\partial_x^2+2b\partial_x\partial_y+c\partial_y^2\) 依照 \(D=b^2-ac\) 的符号分别成为
\[
 \partial_r^2+\partial_s^2\quad(D<0),\quad
 \partial_r^2-\partial_s^2\quad(D>0),\quad
 \partial_r^2\quad(D=0).
\]
低阶项须按同一个变换保留。双曲标准型还可用 \(\xi=r+s\)、\(\eta=r-s\) 化为 \(4\partial_\xi\partial_\eta\)。
\end{theorem}
\begin{proof}
先设 \(a\ne0\)，取 \(r=x\)、\(s=y-(b/a)x\)。此变换行列式为一，且
\[
 \partial_x=\partial_r-\frac ba\partial_s,\quad
 \partial_y=\partial_s.
\]
代入并展开，混合项的系数 \(-2b+2b\) 抵消，得到
\[
 a\partial_r^2+\delta\partial_s^2,\quad
 \delta=c-\frac{b^2}{a}=-\frac Da.
\]
若 \(\delta\ne0\)，再令 \(R=r/\sqrt{|a|}\)、\(S=s/\sqrt{|\delta|}\)，主部变为
\(\operatorname{sgn}(a)\partial_R^2+\operatorname{sgn}(\delta)\partial_S^2\)。乘以 \(\operatorname{sgn}(a)\)，第二项符号等于 \(-\operatorname{sgn}(D)\)，得到前两种形式。若 \(\delta=0\)，只缩放 \(r\) 即得第三种。

若 \(a=0,c\ne0\)，交换 \(x,y\) 后归入上一情形。若 \(a=c=0,b\ne0\)，取 \(r=(x+y)/\sqrt2\)、\(s=(x-y)/\sqrt2\)，则 \(2b\partial_x\partial_y=b(\partial_r^2-\partial_s^2)\)，再缩放并调整符号即可。最后，\(\partial_r=\partial_\xi+\partial_\eta\)、\(\partial_s=\partial_\xi-\partial_\eta\)，相减平方得到 \(4\partial_\xi\partial_\eta\)。
\end{proof}

例如主部 \(4u_{xx}+4u_{xy}+2u_{yy}\) 取 \(r=x/2\)、\(s=y-x/2\) 后成为 \(v_{rr}+v_{ss}\)。同样的变换把
\[
 4u_{xx}+4u_{xy}+u_{yy}+u_x=0
\]
化为 \(v_{rr}+v_r/2-v_s/2=0\)，即 \(v_s=2v_{rr}+v_r\)。二阶部分秩为一，而原来的一阶项恰好提供了新的时间导数。若删掉原式的 \(u_x\)，新方程只是 \(v_{rr}=0\)，便没有这一演化含义。

\ProseParagraphBreak

上述论证适用于常系数。若每一点分别对 \(A(x,y)\) 选一个对角化矩阵，并不能立即把这些矩阵拼成坐标变换的导数：一个候选矩阵场必须满足混合偏导的相容条件。特别是把位置依赖的向量场平方，会额外产生系数导数；不能把变系数算子当成常系数多项式直接配方而忽略这些项。

\begin{proposition}[候选坐标导数的相容性]\label{b2:prop:08-coordinate-integrability}
设 \(n\) 为正整数，\(V\subset\mathbb R^n\) 为关于 \(z_0\in V\) 星形的开集，\(B=(B^k_j)\in C^1(V;\mathbb R^{n\times n})\)。存在 \(F\in C^2(V;\mathbb R^n)\) 满足 \(DF=B\)，当且仅当
\[
 \partial_iB^k_j=\partial_jB^k_i\quad(1\leq i,j,k\leq n).
\]
若还要求 \(\det B\ne0\)，所得 \(F\) 在每一点附近是微分同胚，但未必在整个 \(V\) 上单射。
\end{proposition}
\begin{proof}
必要性是各个分量 \(F^k\) 的混合偏导相等。反之，令 \(h=z-z_0\)，定义
\[
 F^k(z)=\int_0^1\sum_jB^k_j(z_0+th)h_j\,\mathrm{d}t.
\]
线段包含于开集，局部共同紧邻域允许积分求导。对 \(z_i\) 求导并用相容条件，得到
\[
 \begin{aligned}
 \partial_iF^k(z)
 &=\int_0^1\left(B^k_i(z_0+th)
       +t\sum_jh_j\partial_iB^k_j(z_0+th)\right)\,\mathrm{d}t\\
 &=\int_0^1\frac{\mathrm{d}}{\mathrm{d}t}
       \left(tB^k_i(z_0+th)\right)\,\mathrm{d}t=B^k_i(z).
 \end{aligned}
\]
由于 \(B\in C^1\)，这个等式使 \(F\in C^2\)。行列式非零时，用逆函数定理得到局部可逆性。
\end{proof}

例如 \(A=\operatorname{diag}(1,e^{-2x})\) 在每一点都可用 \(B=\operatorname{diag}(1,e^x)\) 化成 \(BAB^{\mathsf T}=I\)。但若 \(B=DF\)，第二行会要求 \(F^2_x=0\)、\(F^2_y=e^x\)，其混合偏导分别为 \(0\)、\(e^x\)，矛盾。这只排除这一候选矩阵场，不排除另外求解坐标方程的可能性。即使相容性成立，全局单射还须另核：\(F(x,y)=(e^x\cos y,e^x\sin y)\) 的行列式处处为 \(e^{2x}>0\)，却以 \(2\pi\) 为周期重复覆盖。局部与全域标准化都不能只靠逐点矩阵运算完成。

\subsection{特征坐标}
\label{b2:subsec:080302}

先处理一个双曲邻域。假定 \(U\subset\mathbb R^2\) 开，\(a,b,c\in C^2(U;\mathbb R)\)，且在 \(U\) 上 \(a\ne0\)、\(D=b^2-ac>0\)。二维特征斜率为
\[
 m_\pm=\frac{b\pm\sqrt D}{a}.
\]
两族曲线分别由 \(y'=m_\pm(x,y)\) 给出。以横截直线 \(x=x_0\) 上交点的 \(y\) 坐标为标签，可以沿曲线构造新坐标；为对方程使用二阶链式法则，必须同时证明这些标签为 \(C^2\)。下面把标签的构造和这一步正则性一起写出。

对每一族，令 \(Y_\pm(x;s)\) 解
\[
 \partial_xY_\pm=m_\pm(x,Y_\pm),\quad Y_\pm(x_0;s)=s.
\]
由于 \(m_\pm\in C^2\)，可把 \(x\) 加入状态，令其导数为 \(1\)，再用\cref{b2:thm:07-local-ode,b2:thm:07-flow-derivative}。在 \((x_0,y_0)\) 附近，所有邻近标签 \(s\) 的解因而在共同短区间上存在，并对 \((x,s)\) 为 \(C^1\)。先固定一族，暂写 \(m,Y\)，记 \(P=Y_s\)。一阶变分方程为
\[
 P_x=m_y(x,Y)P,\quad P(x_0;s)=1,
 \quad
 P(x;s)=\exp\!\left(\int_{x_0}^x m_y(r,Y(r;s))\,\mathrm dr\right)>0.
\]
二阶初值导数的候选值 \(Q\) 由线性方程
\[
 Q_x=m_y(x,Y)Q+m_{yy}(x,Y)P^2,
 \quad Q(x_0;s)=0
\]
确定；积分因子给出它的具体表达
\[
 Q(x;s)=P(x;s)\int_{x_0}^x m_{yy}(r,Y(r;s))P(r;s)\,\mathrm dr.
\]
零初值来自 \(Y(x_0;s)=s\)，其对独立标签 \(s\) 的二阶导数为零。

验证 \(Q=Y_{ss}\) 时，在共同短区间内令 \(Y_\varepsilon(x)=Y(x;s+\varepsilon)\)，并对 \(P\) 作差商
\[
 D_\varepsilon(x)=\frac{P(x;s+\varepsilon)-P(x;s)}{\varepsilon}.
\]
从两个一阶变分方程相减，得到
\[
 (D_\varepsilon)_x
 =m_y(x,Y_\varepsilon)D_\varepsilon
  +\frac{m_y(x,Y_\varepsilon)-m_y(x,Y)}{\varepsilon}P,
 \quad D_\varepsilon(x_0)=0.
\]
一阶 \(C^1\) 依赖保证 \((Y_\varepsilon-Y)/\varepsilon\to P\) 在紧 \(x\) 区间上一致成立。中值积分又给出
\[
 \frac{m_y(x,Y_\varepsilon)-m_y(x,Y)}{\varepsilon}
 =\int_0^1m_{yy}\bigl(x,Y+\theta(Y_\varepsilon-Y)\bigr)
       \frac{Y_\varepsilon-Y}{\varepsilon}\,\mathrm d\theta.
\]
缩小邻域，使所有曲线及连接线段留在 \(U\) 内的同一个紧集上。\(m_{yy}\) 在此一致连续，因此该差商一致趋于 \(m_{yy}(x,Y)P\)。将 \(D_\varepsilon\) 与 \(Q\) 的积分方程相减，系数有共同界，非齐次余项一致趋零，Gronwall 估计便给出 \(D_\varepsilon\to Q\) 一致。再对邻近的 \(s\) 比较 \(Q\) 的线性方程，同一估计说明 \(Q\) 对 \((x,s)\) 连续。于是 \(Y_{ss}=Q\) 存在且连续；在 \(x<x_0\) 的区间上反转积分方向即可作同样的估计。

余下的二阶导数由原方程直接得到：
\[
 Y_{xx}=m_x(x,Y)+m_y(x,Y)Y_x,
 \quad Y_{xs}=m_y(x,Y)P.
\]
它们也连续，所以 \(Y\) 对 \((x,s)\) 联合为 \(C^2\)。由 \(P>0\)，映射 \((x,s)\mapsto(x,Y_\pm(x;s))\) 的行列式非零，\cref{b2:thm:01-inverse}给出局部逆 \(s=S_\pm(x,y)\)。原映射为 \(C^2\)，逆微分公式使逆的导数为 \(C^1\)，故 \(S_\pm\in C^2\)。

取 \(\xi=S_+\)、\(\eta=S_-\)，它们分别在两族特征线上保持常值，称为\term[tezheng-zuobiao]{特征坐标}[characteristic coordinates]。沿轨道求导以及对逆坐标求导，分别给出
\[
 \xi_x+m_+\xi_y=0,\quad
 \eta_x+m_-\eta_y=0,\quad
 \xi_y=\frac1{\partial_sY_+}>0,\quad
 \eta_y=\frac1{\partial_sY_-}>0.
\]
后两式中的分母在各自的逆坐标处取值。这保证了两个非零导数，不能任取两个特征方程解便断言它们构成坐标。

\begin{proposition}\label{b2:prop:08-hyperbolic-coordinates}
设 \(U\subset\mathbb R^2\) 为开集，\(a,b,c\in C^2(U;\mathbb R)\)，\(a\ne0\)、\(D=b^2-ac>0\) 在 \(U\) 上成立。设 \(d_1,d_2,k\in C(U;\mathbb R)\)，\(\mathbb K\in\{\mathbb R,\mathbb C\}\)，\(f\in C(U;\mathbb K)\)，并记
\[
 Lu=au_{xx}+2bu_{xy}+cu_{yy}+d_1u_x+d_2u_y+ku.
\]
对任意 \((x_0,y_0)\in U\)，以直线 \(x=x_0\) 上交点的 \(y\) 坐标为标签，分别沿 \(y'=(b\pm\sqrt D)/a\) 构造 \(\xi,\eta\)。存在该点的开邻域 \(W\subset U\)，使 \(F=(\xi,\eta):W\to\widetilde W=F(W)\) 为 \(C^2\) 微分同胚。对 \(u\in C^2(W;\mathbb K)\) 及 \(v=u\circ F^{-1}\)，方程 \(Lu=f\) 等价于
\[
 \gamma(\xi,\eta)v_{\xi\eta}
 +\beta_1(\xi,\eta)v_\xi+\beta_2(\xi,\eta)v_\eta
 +\kappa(\xi,\eta)v=\widetilde f(\xi,\eta),
 \quad \gamma\ne0,
\]
其中 \(\gamma,\beta_1,\beta_2,\kappa\in C(\widetilde W;\mathbb R)\)，\(\widetilde f=f\circ F^{-1}\in C(\widetilde W;\mathbb K)\)，各系数均由下面的完整链式公式确定。
\end{proposition}
\begin{proof}
坐标的行列式为
\[
 \xi_x\eta_y-\xi_y\eta_x
 =(m_--m_+)\xi_y\eta_y\ne0,
\]
因此可用逆函数定理。所有 \(v\) 的导数在 \(F(x,y)\) 处取值时，对 \(i,j\in\{x,y\}\) 有
\[
 \begin{aligned}
 u_i&=v_\xi\xi_i+v_\eta\eta_i,\\
 u_{ij}&=v_{\xi\xi}\xi_i\xi_j
       +v_{\xi\eta}(\xi_i\eta_j+\eta_i\xi_j)
       +v_{\eta\eta}\eta_i\eta_j
       +v_\xi\xi_{ij}+v_\eta\eta_{ij}.
 \end{aligned}
\]
由特征方程，\(p((x,y),\mathrm d\xi)=p((x,y),\mathrm d\eta)=0\)，新主矩阵的两个对角元为零。混合导数系数满足
\[
 \gamma\circ F
 =2\bigl(a\xi_x\eta_x+b(\xi_x\eta_y+\xi_y\eta_x)+c\xi_y\eta_y\bigr).
\]
合同变换保持矩阵可逆，而 \(\left(\begin{smallmatrix}0&\gamma/2\\\gamma/2&0\end{smallmatrix}\right)\) 可逆恰好要求 \(\gamma\ne0\)。按\cref{b2:prop:08-coordinate-change}，其余系数为
\[
 \begin{aligned}
 \beta_1\circ F&=a\xi_{xx}+2b\xi_{xy}+c\xi_{yy}
                    +d_1\xi_x+d_2\xi_y,\\
 \beta_2\circ F&=a\eta_{xx}+2b\eta_{xy}+c\eta_{yy}
                    +d_1\eta_x+d_2\eta_y,\\
 \kappa&=k\circ F^{-1},\quad \widetilde f=f\circ F^{-1}.
 \end{aligned}
\]
这些表达式的右侧均已在原坐标中确定，再通过 \(F^{-1}\) 视为新坐标的函数。坐标为 \(C^2\)，低阶系数与源项连续，故所得系数和右端连续，逐项代入便得到等价方程。
\end{proof}

也可直接计算这一混合系数。由 \(m_++m_-=2b/a\)、\(m_+m_-=c/a\)，代入梯度关系得到
\[
 \gamma\circ F=2\xi_y\eta_y\bigl(am_+m_--b(m_++m_-)+c\bigr)
 =-\frac{4D}{a}\xi_y\eta_y.
\]
因此在给定双曲邻域内它明确非零，可以把整个变换后方程除以 \(\gamma\)，得到
\[
 v_{\xi\eta}+\frac{\beta_1}{\gamma}v_\xi
 +\frac{\beta_2}{\gamma}v_\eta+\frac{\kappa}{\gamma}v
 =\frac{\widetilde f}{\gamma}.
\]
这里全部商都是 \((\xi,\eta)\) 的函数，右端也必须同步相除。接近 \(D=0\) 的退化边界时，这种除法的统一界却可能消失；规范化后的低阶系数可能变大，点态标准型不附带全域估计常数。

两个特征标签可以交换命名。例如在 \(x>0\) 上，方程 \(u_{xx}-x^2u_{yy}=0\) 的两族特征为 \(y\pm x^2/2\) 恒定。取 \(\xi=y+x^2/2\)、\(\eta=y-x^2/2\)，其行列式为 \(2x\ne0\)。两次求导得到
\[
 u_{xx}-x^2u_{yy}=-4x^2v_{\xi\eta}+v_\xi-v_\eta.
\]
新方程不只是 \(v_{\xi\eta}=0\)；两个一阶项正是坐标曲率带来的贡献。在 \(x=0\) 上行列式退化，所以不能用同一坐标穿过那里。

\ProseParagraphBreak

另设 \(U\subset\mathbb R^2\) 为开集，\(a,b,c\in C^2(U;\mathbb R)\)，\(a\) 处处非零，且 \(D=b^2-ac\equiv0\)。低阶系数取 \(d_1,d_2,k\in C(U;\mathbb R)\)，源项取 \(f\in C(U;\mathbb K)\)，其中 \(\mathbb K\in\{\mathbb R,\mathbb C\}\)，仍对 \(C^2\) 的 \(\mathbb K\) 值函数考虑同样形式的 \(Lu=f\)。令 \(\beta=b/a\)、\(V=\partial_x+\beta\partial_y\)，则
\[
 a\partial_x^2+2b\partial_x\partial_y+c\partial_y^2
 =aV^2-a(\partial_x\beta+\beta\partial_y\beta)\partial_y.
\]
取 \(\xi=x\)，并沿 \(V\) 的流构造 \(\eta\) 使 \(V\eta=0\)、\(\eta_y\ne0\)，则 \(V=\partial_\xi\)。方程成为
\[
 a\,v_{\xi\xi}+\beta_1v_\xi+\beta_2v_\eta+\kappa v=\widetilde f.
\]
这是邻域中的秩一标准形式。只有 \(\beta_2\ne0\) 时才能解出 \(v_\eta\)，要按 \(\eta\) 正方向解释为向前热型，还须 \(-a/\beta_2>0\)。秩为一的二阶部分本身没有决定这些条件。

\ProseParagraphBreak

例如 \(u_{xx}-yu_y=0\) 在整个平面上有同一个秩一主矩阵。在 \(y>0\) 上，方程可写成 \(u_y=y^{-1}u_{xx}\)，沿增大的 \(y\) 是正扩散；在 \(y<0\) 上同一方向却给出负扩散；在 \(y=0\) 上，法向一阶系数也消失，不能用相同除法解出 \(u_y\)。类型名称没有变化，演化解释的符号和退化位置却都发生了变化。低阶项在这种问题中承担了最高时间导数的职责。

\subsection{椭圆主部的特殊坐标与极坐标}
\label{b2:subsec:08-elliptic-polar-coordinates}

椭圆型没有两族实特征线，因此上述沿实轨道找坐标的方法不适用。常系数主部及一点处冻结系数的主部可以用上述线性变换标准化；一般变系数椭圆主部的邻域标准化还需解相应的坐标方程。某些特殊系数可直接处理，例如在全平面上取 \(r=e^x\)、\(s=y\)，有
\[
 u_{xx}+e^{2x}u_{yy}=r^2(v_{rr}+v_{ss})+r v_r.
\]
这个映射是全平面到半平面的微分同胚。除以 \(r^2\) 后主部化为 Laplace 算子，但保留了 \(r^{-1}v_r\)；所获表达式的定义域也随坐标改变。

\ProseParagraphBreak

极坐标是另一个必须保留低阶项的例子。在不穿过原点且选定角度分支的区域中，令 \(x=r\cos\theta\)、\(y=r\sin\theta\)，\(u(x,y)=v(r,\theta)\)。逆坐标导数给出
\[
 \partial_x=\cos\theta\,\partial_r-\frac{\sin\theta}r\partial_\theta,
 \quad\partial_y=\sin\theta\,\partial_r+\frac{\cos\theta}r\partial_\theta.
\]
平方时必须让导数也作用于系数。例如
\[
 \partial_x\cos\theta=\frac{\sin^2\theta}r,
 \quad\partial_x\left(-\frac{\sin\theta}r\right)
 =\frac{2\sin\theta\cos\theta}{r^2}.
\]
所以
\[
 \begin{aligned}
 u_{xx}
 &=\cos^2\theta\,v_{rr}-\frac{2\sin\theta\cos\theta}r v_{r\theta}
   +\frac{\sin^2\theta}{r^2}v_{\theta\theta}\\
 &\quad+\frac{\sin^2\theta}r v_r
   +\frac{2\sin\theta\cos\theta}{r^2}v_\theta.
 \end{aligned}
\]
\(u_{yy}\) 的表达式可由上式交换 \(\sin^2\theta\) 与 \(\cos^2\theta\)，并将混合导数项与 \(v_\theta\) 项变号得到。相加得到
\begin{equation}\label{b2:eq:08-polar-laplacian}
 \Delta u=v_{rr}+r^{-1}v_r+r^{-2}v_{\theta\theta}.
\end{equation}
若只把常系数方向导数的平方机械相加，就会漏掉 \(r^{-1}v_r\)。

\ProseParagraphBreak

在环域内寻找径向调和函数，方程为 \((rv_r)'=0\)，所以 \(v=C\log r+D\)。若取角模态 \(v=R(r)\cos(k\theta)\)、\(k\geq1\) 为整数，则得到 \(r^2R''+rR'-k^2R=0\)。令 \(s=\log r\)、\(W(s)=R(e^s)\)，有 \(W''-k^2W=0\)，故 \(R=Ar^k+Br^{-k}\)。这些模式在环域都合法；若要求连续延伸到原点，负幂和对数项就必须排除。这里角变量的周期为 \(2\pi\)，因此角模态写成 \(\cos(k\theta)\)，与单位周期变量下写 \(e^{2\pi\mathrm{i}kx}\) 的约定相容。坐标奇性、函数奇性和频率归一化须分别核对。

\subsection{高维主符号}
\label{b2:subsec:080303}

在 \(n\) 维，实对称矩阵可经正交变换化为 \(\operatorname{diag}(\lambda_1,\ldots,\lambda_n)\)，再沿非零特征方向缩放，使非零对角元都成为 \(1\) 或 \(-1\)。因此常系数主部完全由惯性三元组 \((n_+,n_-,n_0)\) 决定，允许整体乘以负数时还可交换前两项。这个结论包含二维的三种标准型，但高维中还有更多可能。

\ProseParagraphBreak

例如 \(\partial_1^2+\partial_2^2-\partial_3^2-\partial_4^2\) 的惯性为 \((2,2,0)\)。它既不椭圆，也不属于只有一个异号方向的双曲类型，更不是抛物型。本章把非退化且正、负惯性指数均至少为 \(2\) 的主部称为\term[chaoshuangquxing]{超双曲型}[ultrahyperbolic]，以区别只有一个异号方向的双曲型；这一范围可对照 Walter A. Strauss 的《Partial Differential Equations: An Introduction》\cite[Section 1.6, pp. 30--31]{Strauss2008Partial}。又如三变量算子 \(\partial_1^2\) 的惯性为 \((1,0,2)\)，有两个完全缺失的二阶方向。对这些算子只说“判别式非正”没有明确含义；一个行列式也不能恢复全部惯性信息。

\ProseParagraphBreak

变系数时，可以在任意固定点 \(z_0\) 选一个线性坐标使 \(A(z_0)\) 标准化。邻近点的矩阵通常只是接近标准矩阵。非退化点附近的类型稳定，可由下面的定量矩阵论证得到。

\begin{lemma}[惯性的小扰动稳定性]\label{b2:lem:08-inertia-perturbation}
设 \(n\) 为正整数，\(A_0\in\mathbb R^{n\times n}\) 为可逆实对称矩阵，\(\mu>0\) 为其各特征值绝对值的最小值，\(\|\cdot\|\) 为欧氏算子范数。若实对称矩阵 \(H\in\mathbb R^{n\times n}\) 满足 \(\|H\|<\mu\)，则 \(A_0+H\) 仍可逆，并具有与 \(A_0\) 相同的惯性。
\end{lemma}
\begin{proof}
令 \(E_+,E_-\) 分别为 \(A_0\) 的正、负特征子空间。对 \(v\in E_+\)，有
\[
 v^{\mathsf T}(A_0+H)v\geq(\mu-\|H\|)|v|^2>0\quad(v\ne0),
\]
所以新矩阵至少有 \(\dim E_+\) 个正特征值。这里使用\secref{b2:sec:0801}中正定子空间最大维数的刻画。对 \(v\in E_-\) 同样有 \(v^{\mathsf T}(A_0+H)v\leq-(\mu-\|H\|)|v|^2\)，故至少有 \(\dim E_-\) 个负特征值。由于 \(A_0\) 可逆，这两个下界的和已为总维数，因此新矩阵没有零特征值，正、负个数也不能再增加，惯性相同。
\end{proof}

连续系数在 \(z_0\) 附近使 \(\|A(z)-A(z_0)\|<\mu/2\)，故类型稳定。对于正定矩阵，以上证明还直接给出最小特征值的局部下界 \(\mu/2\)。秩退化情形没有这样的稳定性：\(\operatorname{diag}(1,\varepsilon)\) 在任意小的正、负 \(\varepsilon\) 下分别椭圆和双曲，而 \(\varepsilon=0\) 时秩为一。稳定结论依赖于非零谱与零之间的实际距离。

\subsection{几何意义}
\label{b2:subsec:080304}

主符号作用在余向量上，特征曲线却画在自变量空间里。两者之间的联系可以用一个直接计算辨明。设 \(n\) 为正整数，\(U\subset\mathbb R^n\) 开，\(\phi\in C^1(U;\mathbb R)\)，在 \(z\in U\) 有 \(\nu=\mathrm d\phi(z)\ne0\)。在欧氏坐标中，把余向量的分量写成列，它对切向量 \(w\) 的作用为 \(\nu^{\mathsf T}w\)；等值面 \(\phi=\phi(z)\) 在该点的切空间正是 \(\ker\nu\)。若实对称主矩阵 \(A(z)\) 满足 \(p(z,\nu)=0\)，则
\[
 \nu^{\mathsf T}A(z)\nu=0.
\]
所以当 \(A\nu\ne0\) 时，\(A\nu\) 是等值面内的一个非零切向量。对于非退化双曲型，\(A\nu\) 必非零；对于秩退化类型，它可能恰为零。余向量 \(\nu\) 与向量 \(A\nu\) 的角色不同：前者测试某个方向是否切于等值面，后者是在这个测试下取零的切向量。

\ProseParagraphBreak

例如对 \(c>0\) 的 \(1+1\) 维波动方程 \(u_{tt}-c^2u_{xx}=0\)，固定自变量顺序为 \(z=(t,x)\)，则
\[
 A=\begin{pmatrix}1&0\\0&-c^2\end{pmatrix},\quad
 \phi(t,x)=x-ct,\quad \nu=\mathrm d\phi=(-c,1)^{\mathsf T}.
\]
直接计算得到
\[
 p(z,\nu)=(-c)^2-c^2=0,\quad
 A\nu=(-c,-c^2)^{\mathsf T}=-c(1,c)^{\mathsf T}.
\]
对固定 \(q\in\mathbb R\)，直线 \(x-ct=q\) 可参数化为 \((t,x)=(r,cr+q)\)；其切向量为 \((1,c)^{\mathsf T}\)，确与 \(A\nu\) 平行。\(\nu\) 记录的是余空间中的特征法向，\(A\nu\) 则给出这条时空直线的切向，二者不能互相替代。

\ProseParagraphBreak

在一维波动方程中，特征曲线为 \(x\pm ct\) 恒定，两族曲线把初值上的区间与后来一点联系起来，\chapref{b2:ch:09}将据此推导表示公式。高维波动方程在余空间中对应实特征锥，\chapref{b2:ch:11}进一步研究它与传播区域的关系。热方程的时间方向由一阶项选定，不能把秩退化的零方向画成波动锥；Laplace 方程则没有非零实特征余向量。坐标标准化明确了这些差别，具体解的性质仍由后续各章在相应模型中建立。


\begin{chaptersummary}
主符号把最高阶导数的系数组织为法向上的二次型。坐标变换使主矩阵作合同变换，因此秩和惯性保持不变；具体特征值、统一常数和低阶项则需要另行追踪。二维非零实主部的三类由判别式控制，高维还存在多个正负方向以及更高维的退化核。

\ProseParagraphBreak

双曲特征坐标来自两族不同的实轨道，秩一坐标来自一族轨道；椭圆情形没有同样的实特征构造。某些特殊椭圆系数可以直接换元，极坐标也能适应径向与角模态，但坐标产生的一阶项、角度分支及定义域都必须保留。常系数标准化不能直接解释成任意变系数方程的邻域标准化。尤其在热型问题中，时间方向由一阶项和空间二阶项共同确定。分类整理了方程的结构，存在唯一性、稳定性和传播仍要对各模型分别建立。
\end{chaptersummary}

\BookExercises

下面各题把分类与坐标计算、退化、数据稳定性相结合，最后一题从原方程构造两族特征坐标，再把原坐标中的数据转为标准型里的自由函数。题目中的变换均须核对可逆性；需要说明“不可能”时，应给出保持不变的量或一个具体反例。

% 改编：Strauss2008Partial Section1.6, Exercise3，官方样章印刷31页；从二维旋转不变推广到n维全部正交变换，并展开低阶系数证明。
\Needspace{6\baselineskip}
\begin{exercise}\label{b2:ex:08-orthogonal-invariance}
设 \(n\geq2\) 为整数，\(L=\sum_{i,j=1}^na^{ij}\partial_{ij}+\sum_{i=1}^n b^i\partial_i+c\) 为作用在 \(\mathbb R^n\) 上的实常系数算子，其中 \(A=(a^{ij})\in\mathbb R^{n\times n}\) 对称，\(b\in\mathbb R^n\)、\(c\in\mathbb R\)。证明 \(L\) 在所有正交坐标变换下不变，当且仅当 \(L=\alpha\Delta+c\)，其中 \(\alpha\in\mathbb R\)。
\end{exercise}
\begin{solution}
正交变换 \(y=Qz\) 把主矩阵变为 \(QAQ^{\mathsf T}\)，一阶系数变为 \(Qb\)。取 \(Q=-I\)，不变性迫使 \(b=-b\)，故 \(b=0\)。再取只反射一个坐标的对角正交矩阵，所有连接该坐标与其他坐标的非对角元变号，故这些元均为零。最后取交换两个坐标的正交矩阵，各对角元必须相等，得到 \(A=\alpha I\)。反过来，\(Q\alpha I Q^{\mathsf T}=\alpha I\)，零阶项也不变，故这类算子具有所需不变性。
\end{solution}

% 自拟：坐标曲率产生时间导数，而非把低阶项当作无关余项。
\Needspace{6\baselineskip}
\begin{exercise}\label{b2:ex:08-parabolic-curvature}
对 \(u\in C^2(\mathbb R^2;\mathbb C)\)，证明方程
\[
 u_{xx}+2x u_{xy}+x^2u_{yy}=0
\]
处处具有秩一主部，并用 \(\xi=x\)、\(\eta=y-x^2/2\) 将其化为向前热方程。解释为何它没有显式一阶项，却能在新坐标中出现时间导数。
\end{exercise}
\begin{solution}
主矩阵为 \((1,x)^{\mathsf T}(1,x)\)，秩恒为一。给定变换的行列式为一，逆为 \(x=\xi\)、\(y=\eta+\xi^2/2\)。令 \(V=\partial_x+x\partial_y\)，则 \(V=\partial_\xi\)，而 \(V^2u=u_{xx}+2x u_{xy}+x^2u_{yy}+u_y\)。故原式化为 \(v_{\xi\xi}-v_\eta=0\)。时间导数来自坐标的二阶导数，等价地来自求导作用于向量场系数 \(x\)；若只作形式平方而漏掉这一项，便会得到错误方程。
\end{solution}

% 自拟：类型变化与一致常数的定量退化。
\Needspace{6\baselineskip}
\begin{exercise}\label{b2:ex:08-changing-type}
对全平面上的实系数算子 \(L=\partial_x^2+2x\partial_x\partial_y+\partial_y^2\) 的主部按位置分类。在条带 \(\{(x,y)\in\mathbb R^2:|x|<1\}\) 上求每一点的单位向量二次型的最小值和最大值，说明它为何在整个条带上不一致椭圆。证明任何包含 \(x=1\) 上某点的开邻域都不能经 \(C^2\) 微分同胚化为固定的非退化标准型。
\end{exercise}
\begin{solution}
矩阵为 \(\left(\begin{smallmatrix}1&x\\x&1\end{smallmatrix}\right)\)，特征向量为 \((1,1)\)、\((1,-1)\)，特征值为 \(1+x,1-x\)。所以 \(|x|<1\) 时椭圆，\(|x|>1\) 时双曲，\(|x|=1\) 时秩一。条带上的单位球最小、最大值分别为 \(1-|x|,1+|x|\)，前者没有统一正下界。任何穿过 \(x=1\) 的邻域既含秩一的点，也含非退化的点；合同变换保持秩，故不可能得到处处非退化的固定标准型。
\end{solution}

% 自拟：在曲线数据上重建混合导数方程的两项自由函数。
\Needspace{6\baselineskip}
\begin{exercise}\label{b2:ex:08-curved-cauchy}
设 \(I\subset\mathbb R\) 为开区间，\(h\in C^2(I;\mathbb R)\)、\(g\in C^2(I;\mathbb C)\)、\(k\in C^1(I;\mathbb C)\)，且 \(h'(\xi)\ne0\) 对所有 \(\xi\in I\) 成立。对复值方程 \(u_{\xi\eta}=0\)，在曲线 \(\eta=h(\xi)\)、\(\xi\in I\) 上指定
\[
 u(\xi,h(\xi))=g(\xi),\quad
 u_\eta(\xi,h(\xi))=k(\xi).
\]
在曲线上任一点附近构造唯一 \(C^2\) 解，并说明 \(h'\ne0\) 与非特征条件的关系。
\end{exercise}
\begin{solution}
在小矩形中积分 \(u_{\xi\eta}=0\)，先得 \(u_\eta\) 只依赖 \(\eta\)，再积分得 \(u=F(\xi)+G(\eta)\)。因为 \(h'\ne0\)，可局部反解 \(\xi=h^{-1}(\eta)\)。取一个原函数满足 \(G'(\eta)=k(h^{-1}(\eta))\)，再令 \(F(\xi)=g(\xi)-G(h(\xi))\)，得到所求 \(C^2\) 解。改变 \(G\) 的积分常数只会使 \(F\) 反向改变同一常数，故 \(u\) 唯一。主符号为 \(p(\nu)=\nu_\xi\nu_\eta\)，曲线法向为 \((-h',1)\)，所以非特征条件恰为 \(h'\ne0\)。
\end{solution}

% 自拟：高维惯性(2,2,0)排除所有实根时间方向。
\Needspace{6\baselineskip}
\begin{exercise}\label{b2:ex:08-no-time-ultra}
对 \(\xi\in\mathbb R^4\) 令 \(p(\xi)=\xi_1^2+\xi_2^2-\xi_3^2-\xi_4^2\)。证明不存在满足 \(p(\nu)\ne0\) 的实方向 \(\nu\in\mathbb R^4\)，使对每个 \(\eta\in\mathbb R^4\)，二次多项式 \(p(\eta+s\nu)\) 的两个复根（按重数计）均为实数。
\end{exercise}
\begin{solution}
记对称双线性型为 \(B(v,w)=v_1w_1+v_2w_2-v_3w_3-v_4w_4\)。若 \(p(\nu)>0\)，在正子空间 \(P=\operatorname{span}(e_1,e_2)\) 上，线性条件 \(B(\eta,\nu)=0\) 至少有一个非零解 \(\eta\)。于是 \(p(\eta)>0\)，并且
\(p(\eta+s\nu)=p(\eta)+s^2p(\nu)>0\)，没有实根。若 \(p(\nu)<0\)，改在负子空间 \(\operatorname{span}(e_3,e_4)\) 中取 \(\eta\)，所得多项式恒负。两种情形都与要求矛盾。
\end{solution}

% 自拟：椭圆非特征面不蕴含Cauchy稳定，以精确解构造。
\Needspace{6\baselineskip}
\begin{exercise}\label{b2:ex:08-elliptic-cauchy-instability}
圆周 \(\mathbb T=\mathbb R/\mathbb Z\) 取总质量为 \(1\) 的 Lebesgue 测度。在 \(\mathbb T\times[0,\infty)\) 上，对每个正整数 \(n\) 令
\[
 u_n(x,y)=\frac{e^{-\sqrt n}}{2\pi n}\sin(2\pi n x)\sinh(2\pi n y).
\]
证明这些实值光滑函数满足 \(u_{n,xx}+u_{n,yy}=0\)，且对每个固定的非负整数 \(m\)，
\[
 \|u_n(\cdot,0)\|_{C^m(\mathbb T)}
 +\|\partial_yu_n(\cdot,0)\|_{C^m(\mathbb T)}\longrightarrow0,
\]
但对任意固定 \(y_0>0\)，有 \(\|u_n(\cdot,y_0)\|_{L^2(\mathbb T)}\to\infty\)。这里
\(\|v\|_{C^m(\mathbb T)}=\max_{0\leq j\leq m}\|\partial_x^jv\|_{L^\infty(\mathbb T)}\)；若用区域 \(y>0\) 的外法向导数替代 \(\partial_y\)，结论不变。
\end{exercise}
\begin{solution}
对 \(x,y\) 各求两次导数，两个 \((2\pi n)^2\) 项符号相反，故 \(u_{n,xx}+u_{n,yy}=0\)。边界值恒为零，正 \(y\) 导数为 \(e^{-\sqrt n}\sin(2\pi n x)\)，外法向导数只相差负号。因为 \(2\pi n\geq1\)，对每个 \(0\leq j\leq m\)，其 \(j\) 阶切向导数一致范数都不超过 \((2\pi n)^m e^{-\sqrt n}\to0\)，这就同时控制了所规定的整个 \(C^m\) 范数。另一方面，
\[
 \|u_n(\cdot,y_0)\|_{L^2(\mathbb T)}
 =\frac{e^{-\sqrt n}}{2\sqrt2\pi n}\sinh(2\pi n y_0)\longrightarrow\infty.
\]
初始法向 \((0,1)\) 是非特征的，因为其主符号值为一。这一代数条件并未阻止高频数据向内部指数放大。
\end{solution}

% 自拟：拟线性方程在给定梯度上的分类，并比较两种等价表达的统一常数。
\Needspace{6\baselineskip}
\begin{exercise}\label{b2:ex:08-quasilinear-ellipticity}
设 \(U\subseteq\mathbb R^2\) 为开集，\(u\in C^2(U;\mathbb R)\)。对方程
\[
 (1+u_y^2)u_{xx}-2u_xu_yu_{xy}+(1+u_x^2)u_{yy}=0,
\]
证明给定梯度 \(\nabla u\) 所确定的主矩阵正定，求其特征值。比较原式与除以 \((1+|\nabla u|^2)^{3/2}\) 后的方程，在 \(U\) 上有 \(|\nabla u|\leq M\)、\(M\geq0\)，以及 \(|\nabla u|\) 在 \(U\) 上无界这两种情形的一致椭圆性。
\end{exercise}
\begin{solution}
令 \(q=(u_x,u_y)^{\mathsf T}\)，主矩阵为 \(A=(1+|q|^2)I-qq^{\mathsf T}\)。若 \(q\ne0\)，沿 \(q\) 的特征值为一，垂直方向的特征值为 \(1+|q|^2\)；若 \(q=0\)，两个特征值均为一。因此处处正定。梯度有界时，原矩阵的上下界为 \(1,1+M^2\)；相除后的两个特征值为 \((1+|q|^2)^{-3/2}\) 与 \((1+|q|^2)^{-1/2}\)，也有统一正下界和有限上界。若梯度在区域内无界，原矩阵失去统一上界，相除后的矩阵则失去统一正下界。乘一个处处正的函数保持点态类型，但不无条件保持本章同时包含上下界的一致定义。
\end{solution}

% 自拟：从非线性剪切逆推出完整椭圆算子，含必需的一阶补偿。
\Needspace{6\baselineskip}
\begin{exercise}\label{b2:ex:08-nonlinear-shear}
对 \(u\in C^2(\mathbb R^2;\mathbb C)\)，证明全平面方程
\[
 u_{xx}-6x^2u_{xy}+(1+9x^4)u_{yy}-6x u_y=0
\]
经 \(r=x\)、\(s=y+x^3\) 化为 Laplace 方程。核对变换全局可逆，并说明原主矩阵虽然处处椭圆，却在全平面不一致椭圆。
\end{exercise}
\begin{solution}
变换的逆为 \(x=r\)、\(y=s-r^3\)，导数矩阵 \(B=\left(\begin{smallmatrix}1&0\\3x^2&1\end{smallmatrix}\right)\) 行列式为一。原主矩阵为 \(A=B^{-1}B^{-\mathsf T}\)，故 \(BAB^{\mathsf T}=I\)。非线性坐标的二阶导数只有 \(s_{xx}=6x\)，产生一阶项 \(6xv_s\)；原来的 \(-6xu_y=-6xv_s\) 恰好抵消。因此完整方程为 \(v_{rr}+v_{ss}=0\)。原二次型为 \((\xi-3x^2\eta)^2+\eta^2\)，所以正定。其行列式恒为一而迹为 \(2+9x^4\)，最大特征值无界，最小特征值因此趋零。这个坐标变换的导数及逆导数均无全域统一界，与统一常数的失败一致。
\end{solution}

% 自拟：从混合主部求两族特征坐标，再把原坐标的法向数据转换为两自由函数；不以第十一章波动表示公式为前置。
\Needspace{6\baselineskip}
\begin{exercise}\label{b2:ex:08-characteristic-cauchy-data}
设 \(g\in C^2(\mathbb R;\mathbb C)\)、\(h\in C^1(\mathbb R;\mathbb C)\)。对全平面方程
\[
 u_{xx}-2u_{xy}-3u_{yy}=0,
\]
计算判别式与两族特征直线，构造可逆的实线性坐标 \((\xi,\eta)\)，并化为混合导数标准型。进而求满足
\[
 u(0,y)=g(y),\qquad u_x(0,y)=h(y)\quad(y\in\mathbb R)
\]
的唯一 \(u\in C^2(\mathbb R^2;\mathbb C)\)，说明这些原坐标数据如何确定标准型中的两项自由函数。最后以 \(g(y)=y^2\)、\(h(y)=0\) 核对所得公式。
\end{exercise}
\begin{solution}
这里 \(a=1,b=-1,c=-3\)，故 \(D=b^2-ac=4>0\)。特征斜率由 \(m^2+2m-3=0\) 给出，即 \(m=1,-3\)。两族直线分别为 \(y-x\) 恒定和 \(y+3x\) 恒定。取
\[
 \xi=y-x,\quad\eta=y+3x,
 \qquad x=\frac{\eta-\xi}{4},\quad y=\frac{3\xi+\eta}{4}.
\]
正向变换的行列式为 \(-4\)，因此它是全平面的线性微分同胚。令 \(u(x,y)=v(\xi,\eta)\)，则
\[
 \partial_x=-\partial_\xi+3\partial_\eta,\qquad
 \partial_y=\partial_\xi+\partial_\eta,
\]
展开得到
\[
 \partial_x^2-2\partial_x\partial_y-3\partial_y^2
 =-16\partial_\xi\partial_\eta.
\]
故 \(v_{\xi\eta}=0\)。对固定 \(\eta\)，\(v_\eta\) 在整条 \(\xi\) 直线上为常数；再对 \(\eta\) 积分，得到全部 \(C^2\) 解 \(v(\xi,\eta)=F(\xi)+G(\eta)\)。

在 \(x=0\) 上有 \(\xi=\eta=y\)。初值和原法向导数因此要求
\[
 F(y)+G(y)=g(y),\qquad -F'(y)+3G'(y)=h(y).
\]
对第一式求导并联立第二式，得到
\[
 F'=\frac{3g'-h}{4},\qquad G'=\frac{g'+h}{4}.
\]
令 \(H(s)=\int_0^s h(r)\,\mathrm dr\)，可取
\(F=(3g-H)/4\)、\(G=(g+H)/4\)。由给定正则性，二者均为 \(C^2\)，所以所求解为
\[
 u(x,y)=\frac{3g(y-x)+g(y+3x)}4
       +\frac14\int_{y-x}^{y+3x}h(s)\,\mathrm ds.
\]
两自由函数的导数已经确定；其积分常数只能相反地改变，二者之和不变，故 \(u\) 唯一。数据线 \(x=0\) 的法向为 \((1,0)\)，其主符号值为 \(1\ne0\)，与这一可唯一恢复的计算相符。代入 \(g(y)=y^2,h=0\)，得到 \(u(x,y)=y^2+3x^2\)；直接计算有 \(u_{xx}=6\)、\(u_{xy}=0\)、\(u_{yy}=2\)，方程和两项数据均满足。
\end{solution}
